% Français 9118, batch 3 — Gallica views 41–60.
% Pass: Fable 5.1 (claude-fable-5-1), 2026-09-10 — first pass, unchecked against
% the views by a human.
%
% What the twenty views hold. Three things, in this order: the end of the
% « Réponse aux questions de Mlle … » begun on v. 40 (a third hand; the
% lever, the infinite, the hyperbola, 0/0, then two bibliographical notes);
% two drafts in Germain's own hand of her letters to Gauss, signed
% « Le Blanc », heavily corrected; and seven pieces from Legendre — a note
% on Euler's § 47 (the elastic rod with a stylus at its midpoint) and six
% letters dated 22 octobre 1811, 10 novembre 1811, 4 décembre 1811,
% 4 décembre 1813, 19 janvier 1814 (postmarked 1811, see below) and
% 31 décembre 1819. All in French; the mathematics set as formulae are
% Germain's (Gauss's art. 357 equation and a Lagrange 1775 expression),
% Lagrange's plate equation as Legendre reports it, Euler's lamina
% coefficients, and one Diophantine remark of 1819.
%
%   v. 41    f. 35v–36r « Réponse aux questions », continued: the infinite
%                       line struck at a point cannot move; the hyperbola
%                       between its asymptotes, with a sketch on f. 36r.
%   v. 42    f. 36v     the same: 0/0 « egal a tout ce qu'on veut »; where
%                       to find d'Alembert's Dynamique (Firmin Didot) and
%                       Bossut's algebra. f. 37r blank.
%   v. 43    f. 38r     Germain to Gauss, draft, unsigned here: the
%                       Disquisitiones, the theorem 4(x^n−1)/(x−1) = y²±nz²
%                       generalised to x^{n^s}, art. 357. f. 37v blank.
%   v. 44    f. 38v–39r the same: Fermat's equation and n = p−1 for p of
%                       the form 8k+7 (one paragraph struck and rewritten);
%                       Lagrange's Berlin 1775 memoir, p. 352, the form
%                       t²−11u²; Y' and Z' from Y and Z.
%   v. 45    f. 39v     end of the draft: address « chez Silvestre de Sacy,
%                       rue Haute-Feuille », signed « Le Blanc ».
%            f. 40r     second draft to Gauss: thanks for the letter and the
%                       memoir delivered late by de Sacy's household; the
%                       Leipzig edition exhausted; ternary forms of
%                       determinant zero.
%   v. 46    f. 40v–41r the same: quaternary forms, D^{n−1} as determinant
%                       of the adjoint, the order n(n−1); the bookseller
%                       Duprat insolvent.
%   v. 47    f. 41v     end of the second draft, three openings tried and
%                       cancelled.
%            f. 42r     Legendre, note on Euler's problem § 47, « simpliciter
%                       fixus », the stylus as a fixed needle.
%   v. 48    f. 42v–43r the same: problems I–VI, the objection, « tout notre
%                       édifice croule »; signed « L. G. ».
%   v. 49    f. 43v     its cover, postmarked janvier 1811.
%            f. 44r     Legendre, Paris 10 novembre 1811: the memoir with
%                       Lacroix, the supplement, the pendulum.
%   v. 51    f. 45v     cover of the preceding.
%            f. 46r     Legendre, Paris 22 octobre 1811: the memoir « n'est
%                       pas perdu », five commissioners named.
%   v. 53    f. 47v     cover of the preceding, postmarked 22 octobre 1811.
%            f. 48r     Legendre, 4 décembre 1811: Lagrange's equation of
%                       the plate, the lamina as a check.
%   v. 54    f. 48v–49r the same: the double integrals, Lagrange p. 148,
%                       Biot's equation, « emporter la palme ».
%   v. 55    f. 49v     cover of the preceding.
%            f. 50r     Legendre, 4 décembre 1813: the analysis not
%                       understood, the four points, dx dy.
%   v. 56    f. 50v–51r the same: « le sort en est jetté », honourable
%                       mention, printing the researches.
%   v. 57    f. 51v     cover of the preceding (name only).
%            f. 52r     Legendre, 19 janvier: the factor sin ½ω, the
%                       coefficients α β γ δ, pages 152–154 of Euler.
%   v. 58    f. 52v–53r the same: γ = γ', the portions LF and EL, sin uω;
%                       signed « L. G. », Paris 19 janvier 1814 (?).
%   v. 59    f. 53v     its cover, postmarked janvier 1811.
%            f. 54r     Legendre, 31 décembre 1819: p^5 = A(p²−q), « voilà
%                       le nœud », the method judged sterile.
%
% Absent from the output: v. 50 (f. 44v–45r, blank, the letter and address
% showing through), v. 52 (f. 46v–47r, same), v. 60 (f. 54v–55r, blank).
% On v. 42 and v. 43 one of the two facing pages is blank. The BnF's
% stamps and the postmarks are not transcribed; a postmark's date is given
% in a note where it can be read, since it dates the undated pieces.
%
% Folios are the pencilled numbers on the rectos (36–55 across the batch,
% 35 having been read on v. 40); a verso is written as the recto's number
% with v.
%
% Dates. Legendre's letter on f. 52–53 is dated « 19 Janvier 1814 » as far
% as the last figure can be read, but its cover (f. 53v) is postmarked
% 1811, its subject is the note on f. 42–43 whose cover is postmarked
% janvier 1811, and it calls that note « ma 1ère note »; the year is set
% with \uncertain and the postmark noted. The two drafts to Gauss carry no
% date; that they answer Gauss's letters of 1805 (batch 2) is the pass's
% inference and stays in a note.
%
% Notation. Germain's « 4(x^n−1)/(x−1) » is set with \frac as she sets the
% fraction. Legendre's straight d in the plate equation (f. 48r) is kept as
% d, unlike Lagrange's round d of v. 39 (batch 2). Euler's coefficients are
% set α, β, γ, δ: Legendre's second letter is a 6-shaped β, and the letter
% read « y » on f. 42r is the γ of f. 52v. The « ζ » of the lamina equation
% is Legendre's 3-shaped zeta. The lever piece writes « quarré », « sçauroit »
% and « asymptotes » as its scribe does. Sections are the pass's own.
\documentclass[11pt,a4paper]{article}
% The preamble shared by every transcription of the mathematicians' manuscripts in Gallica.
%
% It rests on one idea: **uncertainty compounds, it does not resolve.** A
% transcription that smooths over a doubtful word has destroyed the only thing
% it was made for — a reader must be able to tell, at every point, what was on
% the page and what was supplied. The apparatus macros below are therefore the
% critical apparatus, not decoration.
%
% The same commands are recognised by `scripts/render.mjs`, which produces the
% site's reading view, and by `scripts/tei.mjs`. A macro added here must be
% added there too, or rendering fails loudly — which is the wanted behaviour.
%
% Comments here are English, like the rest of the repository. The documents
% this preamble sets are French, and that is deliberate: see the skills.

\ProvidesPackage{germain}[2026/09/15 Transcriptions of mathematicians' manuscripts in Gallica]

\usepackage{amsmath,amssymb,amsthm}
\usepackage{mathrsfs}
% Kept from the parent preamble so the renderer's diagram support stays
% available; these manuscripts draw no commutative diagrams, and nothing here uses it.
\usepackage{tikz-cd}
\usepackage[english,french]{babel}
\usepackage{fontspec}

% --- Greek ------------------------------------------------------------------
% The text face stays fontspec's default, Latin Modern, which has no Greek:
% XeTeX drops every glyph it lacks with a line in the log and nothing on the
% page, and the Greek words the manuscripts quote vanished from the PDFs. So
% the Greek and Greek Extended blocks get a XeTeX character class of their own,
% and entering or leaving a run of it switches the family to CMU Serif — the
% same Computer Modern design, with polytonic Greek — and back. Only the
% family changes: a Greek word in \textit or \textbf keeps its shape and series.
% The transitions to and from every other class carry the tokens that class
% has towards an ordinary letter, so French spacing before « ; » or inside
% « » still holds after a Greek word. No grouping, so no brace or \endgroup
% can come out unbalanced; maths is left alone, since XeTeX inserts nothing
% there. Kept identical in `exercices.sty`.
\makeatletter
\newfontfamily\germainGreekfont{cmun}[
  Extension=.otf, NFSSFamily=germaingreek,
  UprightFont=*rm, ItalicFont=*ti, SlantedFont=*sl,
  BoldFont=*bx, BoldItalicFont=*bi, BoldSlantedFont=*bl]
\newXeTeXintercharclass\germain@greek
\def\germain@greekrange#1#2{%
  \@tempcnta=#1\relax
  \loop
    \XeTeXcharclass\@tempcnta=\germain@greek
    \ifnum\@tempcnta<#2\advance\@tempcnta\@ne\repeat}
\germain@greekrange{"0370}{"03FF}
\germain@greekrange{"1F00}{"1FFF}
\def\germain@greekprev{\rmdefault}
\def\germain@greekin{%
  \edef\germain@greekprev{\f@family}\fontfamily{germaingreek}\selectfont}
\def\germain@greekout{\fontfamily{\germain@greekprev}\selectfont}
\def\germain@greekjoin{%
  \XeTeXinterchartokenstate=\@ne
  \@tempcnta=\z@
  \loop
    \ifnum\@tempcnta=\germain@greek\else
      \XeTeXinterchartoks\germain@greek\@tempcnta=\expandafter{%
        \expandafter\germain@greekout\the\XeTeXinterchartoks\z@\@tempcnta}%
      \XeTeXinterchartoks\@tempcnta\germain@greek=\expandafter{%
        \the\XeTeXinterchartoks\@tempcnta\z@\germain@greekin}%
    \fi
    \ifnum\@tempcnta<4095\advance\@tempcnta\@ne\repeat}
% babel-french sets its own classes, and saves and restores the interchar
% state, each time a language is selected: join again after it does.
\addto\extrasfrench{\germain@greekjoin}
\addto\extrasenglish{\germain@greekjoin}
\AtBeginDocument{\germain@greekjoin}
\makeatother

\usepackage{geometry}
\geometry{a4paper,margin=25mm,marginparwidth=28mm}
\usepackage{marginnote}
\usepackage{xcolor}
\usepackage[normalem]{ulem}
\setlength{\emergencystretch}{3em}
\usepackage{hyperref}
\hypersetup{colorlinks=true,linkcolor=black,urlcolor=black,pdfborder={0 0 0}}

% --- Batch metadata ---------------------------------------------------------
% Taken as they stand by the rendering script, which shows them at the head of
% the reading view. They fix what the file claims to cover. The macro names are
% the parent project's, so that the scripts stay shared; \folder here means the
% volume's slug — `fr-9115` — and \pages counts Gallica views.

\newcommand{\folder}[1]{\gdef\grothFolder{#1}}
\newcommand{\batch}[1]{\gdef\grothBatch{#1}}
\newcommand{\pages}[2]{\gdef\grothFirst{#1}\gdef\grothLast{#2}}
\newcommand{\foldertitle}[1]{\gdef\grothFoldertitle{#1}}

% The holder's own shelfmark, printed as they write it: « Français 9115 ».
\newcommand{\shelfmark}[1]{\gdef\grothShelfmark{#1}}

% The dating, copied verbatim from the holder's catalogue — never inferred
% here. « XIXe siècle » is the BnF's; a pass that dates a leaf from its content
% says so in a \note{}, as its own inference.
\newcommand{\dating}[1]{\gdef\grothDating{#1}}

% The status notice, printed in the title block and repeated as a diagonal
% watermark on every page of the PDF. The manuscripts are in the public
% domain; what this line declares is not a right but a quality — an
% unverified first pass by a machine. The skills write the line; a file
% without it gets no watermark, so leaving it out is visible in the output.
\newcommand{\watermark}[1]{\gdef\grothWatermark{#1}}

\gdef\grothFolder{?}\gdef\grothBatch{}\gdef\grothFirst{?}\gdef\grothLast{?}%
\gdef\grothFoldertitle{}\gdef\grothShelfmark{}\gdef\grothDating{}\gdef\grothWatermark{}

% --- The résumé -------------------------------------------------------------
\newenvironment{resume}
  {\small\setlength{\parskip}{0.45em}}
  {\par\medskip\hrule\medskip}

% --- Keywords ---------------------------------------------------------------
% In English, closing the modernised reading's résumé: the modern vocabulary
% under which to search for what the leaves construct. The manifest extracts
% them to tag the volume — this line is the single source of the tags.
\newcommand{\keywords}[1]{\par\smallskip\noindent
  {\footnotesize\textsc{Keywords} --- \textit{#1}}\par}

% --- The critical apparatus -------------------------------------------------

% The Gallica view number, in the margin. This is the anchor that lets a
% disagreement over a formula be settled by looking at *one* image rather than
% a volume of 750. The site uses it to turn the facsimile as the transcription
% is scrolled. A view is one scanned image: usually one side of a leaf.
\newcommand{\page}[1]{%
  \marginnote{\footnotesize\color{gray!70}\textsf{v.\,#1}}%
  \label{page:#1}\ignorespaces}

% The BnF's pencilled foliation, where the leaf carries it and a pass can read
% it: « 198r ». This is what the literature cites — « ff. 198r–208v » — and
% the one line that turns a citation into a view. Recorded, never inferred:
% a folio a pass cannot read is not written.
\newcommand{\folio}[1]{%
  \marginnote{\footnotesize\color{gray!50}\textsf{f.\,#1}}\ignorespaces}

% The modernised reading's coarser counterpart to \page{}: which views a
% section is drawn from.
\newcommand{\pagerange}[2]{%
  \marginnote{\footnotesize\color{gray!70}\textsf{v.\,#1--#2}}%
  \ignorespaces}

% Illegible: never guessed.
\newcommand{\ill}{\textcolor{red!60!black}{[\,\ldots\,]}}

% A doubtful reading: the word is given, and flagged as uncertain.
\newcommand{\uncertain}[1]{\uline{#1}}

% An editorial addition — a missing letter, an implied word.
\newcommand{\add}[1]{\textcolor{blue!50!black}{[#1]}}

% Struck out by the writer. What was crossed out is part of the document.
\newcommand{\struck}[1]{\sout{#1}}

% The transcriber's note, on the state of the page or the mathematics.
\newcommand{\note}[1]{\par\smallskip{\small\color{gray!60!black}\textit{[#1]}}\par\smallskip}

% A marginal note of hers, distinct from the body of the text.
\newcommand{\marginal}[1]{\par\smallskip\noindent\hspace{1em}%
  {\small\color{gray!60!black}#1}\par\smallskip}

% --- Title ------------------------------------------------------------------
% The title block is French, like the documents themselves.

\AddToHook{shipout/background}{%
  \ifx\grothWatermark\empty\else
    \begin{tikzpicture}[remember picture, overlay]
      \node[rotate=52, scale=3.4, align=center, text=gray!18,
            font=\sffamily\bfseries]
        at (current page.center) {\grothWatermark};
    \end{tikzpicture}%
  \fi}

\AtBeginDocument{%
  \begin{center}
    {\large\bfseries Manuscrits de mathématiciens (Gallica) ---
      \ifx\grothShelfmark\empty\grothFolder\else\grothShelfmark\fi}\\[2pt]
    {\itshape\grothFoldertitle}\\[4pt]
    {\small vues \grothFirst--\grothLast
      \ifx\grothBatch\empty\else\ (lot \grothBatch)\fi}\\[2pt]
    \ifx\grothDating\empty\else
      {\small\color{gray!60!black}Datation du catalogue : \grothDating}\\[2pt]
    \fi
    \ifx\grothWatermark\empty\else
      {\small\bfseries\color{red!45!gray}\grothWatermark}\\[2pt]
    \fi
    {\footnotesize\color{gray!60!black}Transcription établie d'après les images IIIF de Gallica.
     Source gallica.bnf.fr / Bibliothèque nationale de France.\\
     Les lectures incertaines sont soulignées, les illisibles notées [\,\ldots\,],
     les ajouts entre crochets ; \textsf{v.} numérote les vues de Gallica,
     \textsf{f.} la foliotation au crayon de la BnF.}
  \end{center}
  \vspace{1em}}

\endinput


\folder{fr-9118}
\shelfmark{Français 9118}
\batch{3}
\pages{41}{60}
\dating{1801-1900}
\watermark{Lecture automatique\\non vérifiée}
\foldertitle{Correspondance de Sophie GERMAIN avec les mathématiciens et savants Cauchy, Delambre, Fourier, Gauss, Le Gendre, d'Ansse de Villoison, etc.}

\begin{document}

\section{« Réponse aux questions de Mlle », suite : l'infini et l'hyperbole}

\page{41}\folio{35v}
de $S \times CB$. ainsi à une distance 100 fois plus grande il suffira d'une
force cent fois moindre que $P$ pour conserver l'equilibre. et à une distance
infiniment grande une force infiniment petite ou nulle suffira, car le
† \struck{\ill{}}

\marginal{† reciproque de l'infini est 0}

\note{suite du texte de la vue 40 (lot 2), même main ; « ou nulle » est
ajouté dans l'interligne ; le mot biffé après « car le » n'a pu être lu, et le
signe † au-dessus de lui renvoie à la note marginale. Plus haut dans la marge
gauche, au crayon et presque effacé, une formule dont on ne lit que
« $S \times CB$ » au dénominateur, et un mot (« hic » ?) devant une accolade}

ainsi si nous supposons une ligne inflexible et infiniment etendue de $B$ en
$A$, appuyée par son extremité $B$ ou retenue par une puissance quelconque
$R$. une puissance $S$ venant la choquer de $S$ en $C$, ne sçauroit la faire
changer de place, et meme le paradoxe s'evanouira en quelque sorte si l'on
fait attention que pour que la ligne infinie $BA$, prît quelque mouvement, il
faudroit que le point $A$, \struck{parcou} parcourut dans un temps fini un
chemin infiniment grand; ce qui est absurde. mais on tombe dans ces paradoxes
toutes les fois qu'on se transporte dans la region de l'infini; parcequ'il est
absurde qu'on puisse jamais y parvenir. c'est une maniere de parler. dans le
langage rigoureux de la geometrie, l'infini n'est qu'une grandeur plus grande
que toute grandeur donnée.

\note{« inflexible et » est ajouté dans l'interligne au-dessus de « ligne » ;
« $R$ » est écrit au-dessus d'une lettre soulignée deux fois (un $S$ ?) ;
« dans un temps fini » est ajouté dans l'interligne ; « rigoureux » de même}

La Geometrie pure presente dans l'hyperbole un paradoxe semblable. car soit
l'hyperbole entre ses Asymptotes $AC$, \uncertain{$BE$} \&c. dont la nature
est telle

\folio{36r}
que $Cb$, $Cd$, $Ce$ \&c. etant les abscisses et $Bb$, $Dd$, $Ee$ \&c. les
ordonnées, on a toujours le rectangle $CD$ = au quarré $CB$, \struck{$CE$} le
rectangle $CE$ = au meme quarré $CB$ \&c. ainsi à une distance infinie
l'ordonnée $Ee$ sera infiniment petite ou zero. on aura donc le rectangle
d'une ligne infiniment petite ou zero par une ligne infinie egale à un
rectangle donné.

\note{en tête de la page, un croquis : deux asymptotes, l'une verticale
portant $A$ et $C$, l'autre horizontale portant $C$, $b$, $d$, $e$ ; une
hyperbole passant par $B$, $D$, $E$ avec les ordonnées $Bb$, $Dd$, $Ee$
tracées ; une seconde hyperbole plus éloignée marquée $f$, $g$, $h$, coupée
en $i$ par l'horizontale menée de $A$, avec une verticale pointillée
abaissée de $i$}

il y a plus. en supposant qu'on fut parvenu à l'infini on auroit
$0 \times \infty$ egal non seulement au quarré $CB$, mais à toute autre
grandeur finie. car on peut concevoir une autre hyperbole entre les memes
asymptotes, et au point (infiniment distant) ou elle touche \struck{la ligne}
$Cde$ \&c. on aura encore un rectangle infini en longueur, et infiniment
petit ou nul en hauteur qui sera egal a un quarré $Ci$. mais on le repete,
ces paradoxes n'ont lieu que parceque c'est un paradoxe lui meme que de
supposer qu'on puisse atteindre l'infini.

\note{« touche » est ajouté dans l'interligne au-dessus de « la ligne »
biffé}

\page{42}\folio{36v}
La proposition paradoxale que $\frac{0}{0}$ est egal a tout ce qu'on veut
suit en quelque sorte de la precedente. car $\frac{1}{0}$ est $= \infty$ ou a
l'infini. or $0 \times \infty$ peut representer toute grandeur finie, donc a
la place de $\infty$ substituant sa valeur $\frac{1}{0}$ on aura
$0 \times \frac{1}{0}$ ou $\frac{0}{0}$ egal a toute grandeur finie
quelconque mais le premier paradoxe n'ayant lieu qu'a cause d'une supposition
impossible celui ci est dans le meme cas.

\note{un $\infty$ tracé seul dans la marge gauche, en regard de la quatrième
ligne ; un trait horizontal sépare ce paragraphe des deux suivants}

La Dynamique de M. d'Alembert a été jadis imprimée chez Briasson; on
indiquera probablement où elle se trouve aujourd'hui, \struck{\ill{}} chez
\emph{firmin didot}, rue \emph{jadis Dauphine}. ce libraire se chargera meme
probablement de la procurer.

Le traité d'algebre de M. Bossut, est ce que toute reflexion faite, on croit
pouvoir indiquer de mieux pour s'instruire de l'analyse.

\note{une tache couvre le mot avant « été » ; « firmin didot » et « jadis
Dauphine » sont soulignés dans l'original ; un mot biffé avant « chez » n'a
pu être lu. La page en regard (f. 37r) est blanche ; la pièce s'achève ici}

\section{Brouillon d'une lettre de Germain à Gauss, signée « Le Blanc »}

\page{43}\folio{38r}
Vos \emph{disquisitiones} arithmeticæ font depuis long tems l'objet de mon
admiration et de mes études | le dernier chapitre de ce livre renferme entre
autres choses remarquables le beau théorème contenu dans l'équation
$4\frac{(x^n - 1)}{x - 1} = y^2 \pm n z^2$, je crois qu'il peut être
généralisé ainsi $4\frac{(x^{n^s} - 1)}{x - 1} = y^2 \pm n z^2$ $n$ étant
toujours un nombre premier et $s$ un nombre quelconque. Je joins a ma lettre
deux demonstrations de cette généralisation. Après avoir trouvé la premiere
j'ai cherché comment la methode que vous avez emploiée Art. 357 pouvoit etre
appliquée au cas que j'avois a considerer; j'ai fait ce travail avec d'autant
plus de plaisir qu'il m'a fourni l'occasion de me familiariser avec cette
methode, qui je n'en doute pas sera encore dans vos mains l'instrument de
nouvelles decouvertes

\note{main de Germain ; la page en regard (f. 37v) est blanche.
« disquisitiones » est souligné ; un trait vertical après « études » marque
sans doute un alinéa. Aucune date sur le feuillet : ce brouillon paroît être
celui de la première lettre à Gauss, à laquelle répond la lettre de Gauss du
16 juin 1805 (vue 32) — inférence de la présente lecture. L'article 357 est
celui des \emph{Disquisitiones arithmeticae} où l'équation
$4\frac{x^p - 1}{x - 1} = Y^2 \mp pZ^2$ est établie}

\page{44}\folio{38v}
\struck{Je me suis beaucoup occupé de la celebre équation de fermat \ill{}
$x^n + y^n = z^n$ dont l'impossibilité en nombres entiers n'a encore été
demontrée que pour $n = 3$ et $n = 4$ je crois etre parvenu a prouver cette
impossibilité pour $n = p - 1$ $p$ etant un nombre premier de la forme
$8k + 7$.}

\note{ce paragraphe est barré d'un seul trait oblique ; avant
$x^n + y^n = z^n$ un premier essai de l'équation est biffé}

J'ai ajouté a cet Art. quelques autres considerations. la derniere est
relative a la celebre equation de Fermat $x^n + y^n = z^n$ dont
l'impossibilité en nombres entiers n'a encore été demontrée que pour $n = 3$
et $n = 4$ je crois etre parvenu a prouver cette impossibilité pour
$n = p - 1$ $p$ etant un nombre premier de la forme $8k + 7$.

Je prens la liberté de soumettre \struck{\ill{}} ces essais a votre jugement
persuadé que vous ne dedaignerez pas d'éclairer de vos avis un amateur
enthousiaste de la science que vous cultivez avec de si brillans succès.

\note{« ces essais » est écrit dans l'interligne au-dessus de deux ou trois
mots biffés (« mon … travail » ?)}

Rien n'egale l'impatience avec laquelle j'attens la suite du livre que j'ai
entre les mains, je me suis fait informer, que vous y travailliez en ce
moment; \struck{\ill{}} je ne negligerai rien pour me la procurer
\uncertain{aussitôt} qu'elle paroîtra:

\note{un mot biffé après le point-virgule ; « aussitôt » est écrit
par-dessus un premier « plutôt »}

Malheureusement \struck{l'étendue de mes forces} l'etendue de mon esprit ne
repond \struck{guere} pas a la vivacité de mes goûts et je sens qu'il y a une
sorte de \struck{thé} témérité a \struck{\ill{}} importuner un homme de genie
lorsqu'on n'a d'autre titre a son attention qu'une admiration necessairement
partagée par tous ses lecteurs.

\note{« pas » est écrit au-dessus de « guere » biffé ; le mot biffé avant
« importuner » n'a pu être lu}

\folio{39r}
En relisant le memoire de M. de la Grange (Berlin 1775) j'ai vu avec
étonnement qu'il n'a pas su réduire la quantité
\[
s^{10} - 11\left(s^8 - 4 s^6 r^2 + 7 s^4 r^4 - 5 s^2 r^6 + r^8\right) r^2
\quad (\text{p. } 352)
\]
a la forme $t^2 - 11 u^2$; car
\[
s^{10} - 11\left(s^8 - 4 s^6 r^2 + 7 s^4 - 5 s^2 r^6 + r^8\right) r^2
\]
\[
= r^{10} - 2 \cdot 11\, s^6 r^4 + (5 + 6)\, r^8 s^2
- 11\left(s^8 - 6 s^6 r^2 + 9 r^4 s^4 - 2 r^4 s^4\right.
\]

\note{le calcul s'interrompt sur la parenthèse ouverte, un petit trait
au-dessous ; « $7 s^4$ » sans $r^4$ dans la seconde copie, comme écrit ; la
lecture des exposants de la dernière ligne, et celle du premier terme
« $r^{10}$ » (ou $s^{10}$), est douteuse. Le mémoire cité est celui de
Lagrange dans les Mémoires de l'Académie de Berlin pour 1775, la page 352
étant celle du feuillet}

Cette remarque est une nouvelle preuve de l'avantage de votre methode, qui
s'appliquant a toutes les \struck{cas} valeurs de $n$ donne pour chaque cas
des valeurs de $Y$ et $Z$ independantes de tatonnement

si connoissant les valeurs de $Y$ et $Z$ dans l'équation
$4\frac{(x^n - 1)}{x - 1} = Y^2 \pm n Z^2$ on vouloit avoir celles de $Y'$
et $Z'$ dans l'équation $4\frac{(x^n + 1)}{x + 1} = Y'^2 \pm n Z'^2$, il est
clair qu'il suffiroit de changer les signes de tous les termes de $Y$ et $Z$
qui contiennent des puissances \struck{de $x$} dont l'exposant est impair.

Je n'ai pas voulu fatiguer votre attention en multipliant les remarques dont
votre livre a été pour moi l'occasion; si je puis esperer que vous
accueilliez favorablement celles que j'ai l'honneur de

\page{45}\folio{39v}
vous communiquer et que vous ne \struck{dedaigniez} les trouverez pas
entierement indignes de reponse Veuillez l'adresser A Monsieur Silvestre de
Sacy membre de l'institut national. Rue haute feuille A Paris.

Croyez \ill{} au prix que j'attacherois a un mot d'avis de votre part et
agréez l'assurance du profond respect de

Votre très humble serviteur et très assidu lecteur

Le Blanc

\note{« ne » est ajouté dans l'interligne et « dedaigniez » biffé ; après
« Croyez » un blanc de trois ou quatre mots a été laissé, non rempli ; le
feuillet est plus court que les autres. L'adresse est celle que porte la
lettre de Gauss du 16 juin 1805 (vue 35)}

\section{Second brouillon de Germain à Gauss}

\folio{40r}
Je dois vous paroître bien coupable, d'avoir tardé si longtems a vous
remercier de la lettre dont vous m'avez honoré et de l'envoi du memoire que
vous avez bien voulu y joindre cependant il n'y a pas de ma faute le paquet
ne m'a été remis qu'il y a \struck{quatre} huit jours. Mr de Sacy étoit en
voyage depuis plus de deux mois et on avoit négligé chez lui de me faire
tenir. il est vrai que n'esperant pas de vous une reponse si prompte je
n'avois mis aucun soin a m'informer des lettres qui m'etoient adressées.

\note{« huit » est écrit au-dessus de « quatre » biffé. Le mémoire reçu
seroit celui de 1799 que Gauss annonce dans sa lettre du 20 août 1805 (vue
35) — inférence de la présente lecture, le brouillon n'étant pas daté}

Votre mémoire m'a fait d'autant plus de plaisir que je le connoissois déjà
par une lecture rapide que m'avoit procurée l'un des savans auxquels vous
l'avez envoié il y a déjà longtems et qu'ayant toujours eu le desir de
l'étudier comme on doit le faire de tous les ouvrages qui sortent de votre
plume je l'avois inutilement fait demander a Leipsik d'ou j'avois reçu pour
reponse que l'édition étoit epuisée. \struck{Ainsi je vous ai la plus grande
obligation et je \ill{} tout le prix} \struck{Aussi desesperois je de pouvoir
jamais le \uncertain{retrouver}} \struck{et vous ai je la plus grande
obligation de me l'avoir ainsi donné} c'est pourquoi je n'avois pas osé vous
en parler

\note{trois essais de phrase biffés l'un après l'autre ; dans le premier, un
ou deux mots après « et je » n'ont pu être lus}

L'indulgence que vous continuez de me témoigner m'encourage a vous
communiquer encore quelques unes de mes nouvelles recherches, sur
l'arithmetique.

Après avoir reduit suivant que vous l'indiquez, les formes ternaires dont la
déterminante est zéro aux formes binaires j'ai

\page{46}\folio{40v}
cherché si cette propriété ne s'étendoit pas aux formes quaternaires c'est a
dire si ces formes n'etoient pas susceptibles de se réduire aux formes
ternaires lorsque leur déterminante est zéro pour cela j'ai cherché qu'elle
etoit l'expression générale de la déterminante quaternaire, je me suis aidé
pour cela de l'analogie \struck{ce soupçon s'etant verifié j'ai au moyen des
calculs} et j'ai examiné \struck{aussi} ensuite quelques autres proprietés de
ces formes et de leurs adjointes.

\note{la lecture du passage biffé est douteuse ; « ensuite » est écrit
au-dessus de « aussi » biffé}

Je crois qu'en general $D$ etant la déterminante d'une forme composée d'un
nombre $n$ de variables $D^{n-1}$ \struck{sera} est la déterminante de
l'adjointe de cette forme. c'est ainsi que vous avez trouvé $D^2$ pour la
déterminante de l'adjointe ternaire et \struck{vous verrez} que d'après mes
calculs, $D^3$ est la déterminante de l'adjointe quaternaire. \struck{Mais}
cette analogie n'est sans doute pas suffisante pour etablir la generalité de
la proposition. On voit au moins que la déterminante de la forme étant
composée de produits de l'ordre $n$ et les coëfficiens de son adjointe
l'étant de produits de l'ordre $n - 1$, $D^{n-1}$ est du même ordre que la
déterminante de l'adjointe c'est a dire de l'ordre $n(n - 1)$

\note{« est » est écrit au-dessus de « sera » biffé ; « et que » au-dessus
de « vous verrez » biffé ; « etablir la generalité de la proposition » et
« de la forme » sont écrits dans la marge gauche, en regard des lignes où
ils s'insèrent}

\folio{41r}
pour les formes binaires, la déterminante de l'adjointe seroit d'après
\struck{cette proposition la même} chose que celle de la forme. c'est
\struck{ainsi}, que les formes \struck{binaires} binaires n'\struck{auroient}
ces deux propositions savoir que \struck{l'adjointe} la déterminante est de
l'ordre \struck{pour adjointe} $n$ et les coëfficiens de l'adjointe de
l'ordre \uncertain{$0$} — m'ont paru resulter de la nature generale des
formes et de leurs adjointes.

\note{passage très remanié, dont la lecture est incertaine : « ces deux
propositions savoir que … est de l'ordre », « m'ont paru resulter de la
nature generale des formes et » et « de leurs adjointes » sont écrits dans
les interlignes ; « déterminante » est écrit au-dessus du second
« l'adjointe » de la ligne, et l'on ne sait lequel des deux il remplace ; le
chiffre après « de l'ordre » pourroit aussi se lire « $n - 1$ » ; la
première ligne est surmontée d'un trait courbe qui la biffe peut-être. Le
sens visé paroît être que, pour les formes binaires, la déterminante de
l'adjointe est celle de la forme même}

Je regarde comme une faveur la permission que vous voulez bien m'accorder de
vous communiquer mes foibles essais persuadé que vous aurez assez de bonté
pour m'avertir des erreurs qui pourroient m'échapper dans un genre de
recherches \struck{dont} ou vous etes le seul juge éclairé que l'on puisse
consulter.

Les nouveaux renseignemens que j'ai pris au sujet du libraire Duprat ne sont
rien moins que satisfesans, c'est pour cela que j'ai negligé d'en faire le
sujet exprès d'une lettre. Son successeur a dit avoir depuis longtems terminé
ses payemens dont le produit a été aussi tôt dissipé de sorte que votre
debiteur \struck{doit être} \uncertain{sera regardé} \struck{considere} comme
entierement insolvable.

J'aurois desiré \struck{pouvoir} Il est retiré dans une petite ville ou il
vit du \struck{produit} revenu d'un mediocre emploi

\note{« ou » est écrit au-dessus de « dont » biffé ; « revenu » au-dessus de
« produit » biffé ; « J'aurois desiré » n'est pas biffé mais la phrase est
abandonnée. Duprat est le libraire dont Gauss se plaint dans sa lettre du 16
juin 1805 (vue 34)}

\page{47}\folio{41v}
\struck{Je vous prie, Monsieur, de vouloir bien agréer tous mes remercimens
me continuer l'indulgence, que vous avez bien voulu me}

\struck{Il faut bien compter sur votre indulgence pour oser vous adresser un
si volumineux paquet, je sens cependant que votre tems est}

\note{ces deux débuts sont annulés d'un même trait vertical ondulé ; dans le
premier, « tous » est ajouté au-dessus de « mes », et « bienveillance »
au-dessus de « l'indulgence »}

\ill{} l'avis général de toutes les personnes que j'ai consultées
\struck{est} été qu'il est a peu près impossible de tirer de l'argent de lui.

\note{les deux premiers mots de la ligne (« J'a… ») n'ont pu être lus ;
« été » est écrit au-dessus de « est » biffé}

Je n'avois pas jugé necessaire de vous communiquer ces resultats parceque je
ne vois pas que l'on puisse en tirer bon parti et que j'attendois, pour vous
écrire de nouveau que vous m'en eussiez donné la permission, \struck{et} le
retard qu'a occasionné \struck{le voyage} la remise de votre lettre m'a privé
de vous faire plutôt \struck{mes} tous mes remercimens et les protestations
de mon profond respect.

\section{Note de Legendre sur le problème § 47 d'Euler}

\folio{42r}
Euler n'a traité qu'en passant et par forme d'exemple son problème § 47, il
peut s'y être mépris tant en fait de calcul qu'en fait de raisonnement. Il
s'est mépris certainement dans le calcul lorsqu'il a trouvé
$\gamma' = -\gamma$, puisqu'on doit avoir $\gamma' = \gamma$. Il se peut
faire aussi que la seconde solution soit purement analytique et ne
satisfasse pas aux circonstances physiques du problème. C'est ce que je ne
deciderai pas, n'ayant pas assez réfléchi sur ces sortes de questions et
n'ayant pas le loisir ni le goût de me livrer à un examen plus approfondi.
J'aime donc mieux donner cause gagnée à Mlle Sophie que de lutter avec elle
sur un sujet qu'elle a beaucoup médité. Voici seulement ce qui me paroit le
plus probable.

\note{main de Legendre, sans en-tête ni date ; la lettre lue $\gamma$ est
celle de la lettre du 19 janvier (f. 52v), où elle est nettement un gamma.
Le « problème § 47 » est celui du traité d'Euler sur les courbes élastiques
dont la lettre du 19 janvier cite les pages 152 à 157 — identification de
la présente lecture}

Avant toute discussion il faudroit avoir bien fixé le sens des mots
\emph{Simpliciter fixus} qu'emploie Euler. Comme dans ce point $y$ est
toujours zéro, il faut, ce me semble, regarder le stilet comme une aiguille
fixe qui traverse la verge au milieu de sa largeur, et autour duquel elle
peut tourner dans tous les sens — Je ne vois pas quel mot d'Euler puisse
avoir une autre signification.

\note{« Simpliciter fixus » est souligné dans l'original}

\page{48}\folio{42v}
Cela posé si l'on a bien déterminé dans le probl. IV tous les mouvemens que
peut prendre une verge élastique dont les extrémités sont \emph{Simplement
fixes}; parmi tous les mouvemens réguliers possibles il y en aura un certain
nombre dans lesquels le point milieu de la verge demeurera en repos: ces
derniers mouvemens satisferont au probl. du § 47, il ne s'agira donc que de
retrancher de la solution generale du probl. IV, toutes les solutions qui ne
satisfont pas à cette condition.

Pareil raisonnement s'applique à tous les autres cas généraux depuis le
problème I jusqu'au probl. VI, et il s'applique encore au cas où le stilet
seroit appliqué à un autre point que le milieu, ou même aux cas où plusieurs
stilets seroient appliqués en différents points de la verge, au moins
suivant des distances qui seroient dans un rapport rationnel avec la
longueur entiere de la verge.

Cette explication peut faire disparoître beaucoup de difficultés, mais je ne
me dissimule pas qu'elle est sujette à une objection.

Quand on considère dans les problèmes successifs I, II .. VI, les differens
mouvemens de la verge on suppose qu'elle est entierement libre dans les
points intermédiaires, et qu'ils n'éprouvent dans leurs mouvemens aucune
résistance. Le cas n'est plus \struck{\ill{}} le même lorsqu'on conçoit un ou

\note{« Simplement fixes » est souligné ; « le même » est écrit au-dessus
d'un mot biffé}

\folio{43r}
plusieurs stilets appliqués en différens points. Si ces stilets ne
supportent aucune pression dans aucun sens, la solution telle que nous
venons de la concevoir pourra etre appliquée, mais s'ils en supportent une
il faudra y avoir égard, les solutions des problèmes I, II .. VI. ne sont
plus applicables et tout notre édifice croule.

Permettez, mademoiselle, que je vous laisse vous dégager \struck{comme} vous
pourrez de ces ruines, moi je me sauve en vous faisant une très humble
révérence

L. G.

\page{49}\folio{43v}
A Mademoiselle

Mademoiselle Sophie Germain,

Maison de M. Degligny

Rue Ste Croix de la bretonnerie

Paris

\note{adresse au verso du second feuillet de la note ; timbre de la poste
« Janvier … 1811 », le quantième illisible, et une marque « G 5 » ; l'adresse
est écrite en travers}

\section{Lettre de Legendre, Paris, 10 novembre 1811}

\folio{44r}
Mademoiselle, Votre mémoire est en circulation, M. Lacroix l'avoit entre les
mains lundi dernier; je m'informerai demain à qui il l'a remis, et j'y ferai
joindre le Supplément, les commissaires jugeront ensuite s'ils doivent tenir
compte ou non de ce supplément. Je ferai ensorte d'ailleurs que M. de la
Grange ne tarde pas à lire le tout.

Il n'y a pas de difficulté, ce me semble, dans le cas particulier où le
pendule a la vitesse nécessaire pour remonter jusqu'à l'extrémité supérieure
du diametre vertical. Le calcul prouve qu'il faut un tems infini pour que le
pendule arrive à ce point, et alors son mouvement sera anéanti;

Agréez, Mademoiselle, l'hommage de mes sentimens les plus distingués

Le Gendre

Paris ce 10 novembre 1811.

\note{la signature est écrite en deux mots et soulignée, ici et dans les
lettres suivantes}

\page{51}\folio{45v}
A Mademoiselle

Mademoiselle Sophie Germain

Maison de M. Degligny

Rue Ste Croix de la bretonnerie

Paris

\note{adresse de la lettre du 10 novembre 1811, au verso de son second
feuillet ; timbre de la poste « 10 … 1811 » et marque « G 5 »}

\section{Lettre de Legendre, Paris, 22 octobre 1811}

\folio{46r}
Mademoiselle

Votre mémoire n'est pas perdu: il est le seul qu'on ait reçu sur la question
des vibrations des surfaces. on a nommé hier cinq commissaires pour
l'examiner, j'ai l'honneur d'en être un, Mrs Laplace, Lagrange, Lacroix et
Malus sont les quatre autres. Je n'ai rien dit, je vous conseille également
de garder le silence jusqu'au jugement définitif

Je suis avec tous les sentimens que vous me connoissez

Votre dévoué serviteur

Le Gendre

Paris ce 22 8bre 1811

\note{« 8bre » abrège octobre, « bre » en exposant ; un « 7 » au crayon dans
le coin supérieur, à côté du folio}

\page{53}\folio{47v}
Mademoiselle Sophie Germain,

Maison de M. Degligny

Rue St Croix de la bretonnerie

Paris

\note{adresse de la lettre du 22 octobre 1811, au verso de son second
feuillet ; timbre de la poste « Octobre 22 1811 » et marque « G 5 »}

\section{Lettre de Legendre, 4 décembre 1811}

\folio{48r}
Mademoiselle

Je n'ai pas de bonnes nouvelles à vous donner de l'examen du Mémoire. On
trouve que votre équation principale n'est pas exacte, même en admettant
l'hypothèse que l'élasticité en chaque point peut être représentée par
$\frac{1}{r} + \frac{1}{r'}$. M. de la Grange a trouvé que dans cette
hypothèse la vraie équation devroit être de la forme
\[
\frac{d^2 z}{dt^2}
+ k^2\left(\frac{d^4 z}{dx^4} + 2\frac{d^4 z}{dx^2 dy^2}
+ \frac{d^4 z}{dy^4}\right) = 0,
\]
en supposant d'ailleurs $z$ très petit. Je n'ai point vérifié ce calcul; on
peut s'en rapporter à son auteur; mais ce qui du premier coup d'œil confirme
son exactitude, c'est qu'en supposant la surface vibrante réduite à une lame
d'une largeur constante, ce qui peut s'exprimer en faisant
$\frac{dz}{dy} = 0$, on retombe sur l'équation
$\frac{d^2 z}{dt^2} + k^2 \frac{d^4 z}{dx^4} = 0$ qui est, autant qu'il m'en
souvient (car je n'ai pas le volume sous la main) l'équation donnée par Euler
pour les lames elastiques vibrantes. Votre équation ne donneroit pas ce
résultat. La source de votre erreur \struck{\ill{}} paroit être dans la
maniere dont vous avez

\note{Legendre écrit le $d$ droit ; c'est l'équation de la note de Lagrange
copiée sur la vue 39 (lot 2) ; un mot biffé après « erreur »}

\page{54}\folio{48v}
cru pouvoir déduire l'équation de la surface vibrante de l'équation d'une
simple lame; c'est dans les doubles intégrales que vous vous êtes égarée;
Elles ne se prêtent nullement aux substitutions que vous avez employées. Il
falloit pour l'équation de la surface \struck{suivre} la Méthode indiquée
par Lagrange dans sa 1re édition pag. 148, en ajoutant le terme convenable
pour représenter la force due à l'élasticité. Au reste ces choses sont
sujettes à des difficultés particulieres qui n'ont pas encore été bien
éclaircies et il y auroit \struck{même} des objections à faire contre
l'analyse de l'article même que je cite.

M. Biot qui a eu communication de votre Mémoire prétend avoir trouvé la
vraie équation de la surface Elast. vibr., il m'en a communiqué une, qu'il
dit avoir montrée il y a longtems à M. de la Place et qui n'est pas la même
que celle qu'a trouvée M. de la Grange d'après votre hypothèse.

Je n'en rends pas moins justice à des efforts qui sont louables en
eux-mêmes, quoiqu'ils n'ayent pas l'issue que j'aurois desirée, mais c'est
une raison de plus de garder l'incognito

\note{« 1re » avec « re » en exposant : la première édition de la Mécanique
analytique, identification de la présente lecture}

\folio{49r}
et je vous promets de mon côté de garder le plus profond silence.

J'imagine que la même question sera proposée avec un nouveau délai, ainsi
\uncertain{Miséricorde} n'est pas perdue: au contraire il faut plus que
jamais songer à emporter la palme.

Agréez, Mademoiselle, les sentimens affectueux de votre dévoué serviteur

Le Gendre

ce 4 décembre 1811.

\page{55}\folio{49v}
a Mademoiselle

Mademoiselle Sophie Germain,

Maison de M. Degligny,

rue Ste Croix de la bretonnerie

Paris

\note{adresse de la lettre du 4 décembre 1811, au verso de son second
feuillet ; timbre de la poste illisible ; marque « G 5 »}

\section{Lettre de Legendre, 4 décembre 1813}

\folio{50r}
Mademoiselle

Je ne comprends pas du tout l'analyse que vous m'envoyez, il y a certainement
erreur ou dans l'écriture ou dans le raisonnement, et je suis porté à croire
que vous n'avez pas une idée bien nette des opérations qu'on fait sur les
intégrales doubles dans le calcul des Variations. Votre explication des
quatre points ne me satisfait pas davantage. La Grange a eu raison de
considérer deux Elemens consécutifs dans la courbe élastique, et de mesurer
l'élasticité par l'angle compris entre les deux Elemens. On n'a pas
d'Elemens analogues dans les surfaces, ou du moins ceux que vous avez
considérés ne sont pas dans la ligne de l'analogie. Un élément de la surface
a pour projection $dx\,dy$; l'Element suivant a pour projection
$(dx + ddx)(dy + ddy)$; ces projections sont deux carrés séparés; ensuite
l'idée des plans ne s'accommode pas avec ces projections, parcequ'un plan ne
passe pas par quatre points. Il y a donc dans tout cela beaucoup d'obscurité

\page{56}\folio{50v}
Je ne me charge pas de vous lever toutes les difficultés dans une matiere
que je n'ai pas cultivée spécialement et qui n'a pas d'attrait pour moi;
ainsi il est inutile que je vous donne un rendez vous pour en causer.
D'ailleurs le sort en est jetté, il n'y a plus rien à changer au Mémoire, et
avec toute ma bonne volonté je n'y pourrois rien faire.

Il paroit reconnu cependant que votre Equation est réellement celle de la
\struck{surfa} surface vibrante; en mettant l'analyse à part le reste peut
etre bon, en ce qui concerne l'explication des phénomenes. Si la Commission
de l'institut étoit de cet avis, vous pourriez au moins être mentionnée
honorablement, mais je crains bien que l'analyse manquée ne nuise beaucoup
au mémoire, malgré ce qu'il peut contenir de bon.

Dans tous les cas vous aurez la ressource de faire imprimer vos recherches
en rétablissant la vraie analyse ou en la supprimant, et votre travail vous
fera encore honneur. C'étoit peut etre le parti qu'il falloit prendre dans
l'origine,

\folio{51r}
mais je vous promets toujours le plus profond secret, et si vous n'avez pas
commis d'ailleurs quelque indiscretion, la chose sera comme non avenue.

Agréez je vous prie, mes hommages et mon entier dévouement

Le Gendre

ce 4 déc. 1813.

\page{57}\folio{51v}
Mademoiselle

Sophie Germain

\note{adresse de la lettre du 4 décembre 1813, au verso de son second
feuillet, réduite au nom ; ni rue ni timbre}

\section{Lettre de Legendre, Paris, 19 janvier « 1814 » (timbre de 1811)}

\folio{52r}
La Multiplication par $\sin \frac{1}{2}\omega$, contre laquelle je m'étois
élevé dans ma 1ère note, s'explique en examinant les choses de plus près, et
voici comment

Avant de faire aucune supposition sur la valeur de $\omega$, l'auteur (pag.
154) trouve le rapport
\[
\frac{\alpha}{\alpha'}
= \frac{e^{\lambda\omega} - e^{\omega(1 - \lambda)}}
       {e^{\lambda\omega} - e^{-\lambda\omega}},
\]
d'où il conclut $\alpha = e^{\lambda\omega} - e^{\omega(1 - \lambda)}$,
$\alpha' = e^{\lambda\omega} - e^{-\lambda\omega}$; parcequ'en effet il peut
multiplier tous les coefficiens $\alpha, \beta, \gamma, \delta$, par un même
nombre, puisqu'il reste encore un coefficient arbitraire $C$ qui multiplie le
tout.

Mais comme par suite de la valeur $\sin \frac{1}{2}\omega = 0$, on trouve
$\alpha = 0$ et $\alpha' = 0$, il s'ensuit qu'on a mal à propos multiplié
tous les coefficiens $\alpha, \beta, \gamma, \delta$ par une quantité
infinie, puisque le coefficient $\alpha$ qui sans cette multiplication auroit
été zéro, est devenu une quantité finie. Pour rectifier cette erreur il faut
donc supprimer le facteur infini, ou multiplier par $\sin \frac{1}{2}\omega$.

Cette explication laissoit encore quelques obscurités, et il est bien plus
simple de refaire le calcul des coefficiens dans la supposition de
$\sin \frac{1}{2}\omega = 0$ ou $\frac{1}{2}\omega = k\pi$, $\lambda$ étant
$\frac{1}{2}$.

Soit donc $\sin \frac{1}{2}\omega = 0$, et alors en remontant tout
simplement aux équations primitives de la pag. 152, on trouve \struck{\ill{}}
sans aucune difficulté $\alpha = 0$, $\beta = 0$, $\delta = 0$,
$\alpha' = 0$, $\beta' = 0$, $\delta' = 0$, Il ne reste que $\gamma$ et
$\gamma'$ qui ne s'anéantissent pas. Mais les équations dont il s'agit n'en
déterminent pas la valeur et on trouve simplement $\gamma = \gamma'$. A cause

\note{« ma 1ère note » : la note sur le § 47 d'Euler (f. 42–43), dont le
raisonnement est repris ici ; « 1ère » avec « ère » en exposant ; « Multipli
» à la fin de la seconde ligne est raturé et le mot repris ; deux ou trois
symboles biffés après « on trouve » ; « $\delta = 0$ » est écrit
par-dessus une première lettre}

\page{58}\folio{52v}
du Multiplicateur Commun $C$, on peut faire $\gamma = 1$, et on aura
$\gamma' = 1$

Euler dans son analyse (mal ordonnée) trouve $\gamma' = -\gamma$; Mais c'est
une erreur manifeste, et les équations III et IV de la page 152 donnent
évidemment $\gamma = \gamma'$.

Voilà une difficulté très vraie et très grave, et voyez les Conséquences qui
en résultent. Si on a $\gamma = \gamma' = 1$, alors l'équation de la portion
de courbe $LF$ (pag. 157) n'est plus $y = -C \sin(\zeta + \text{\&c})$, mais
bien $y = +C \sin(\zeta + \text{\&c})$ comme celle de la portion $EL$.
\struck{\ill{}}

\note{deux lignes entièrement biffées, d'un trait épais, après « $EL$ » ;
on y distingue « … ces deux portions … seroient … » sans plus}

Il reste donc à chercher laquelle de ces deux équations est la vraie. On
pourroit croire au premier coup d'œil que c'est celle d'Euler, qui semble
indiquer tout de suite des Ordonnées négatives pour la portion $LF$. Hé bien,
point du tout, Euler s'est trompé dans cette équation par suite de son erreur
sur le signe de $\gamma'$, et la vraie équation de la portion $LF$, est
\[
y = C \sin\left(\zeta + t\,\frac{\omega c \sqrt{2gb}}{aa}\right) \cdot
\sin u\omega
\]
absolument comme celle de la portion $EL$, c'est-à-dire que ces deux
portions ne font qu'une seule et même courbe désignée par la même équation.
Résultat qui se rapproche \struck{entierement} de la théorie que Mlle Sophie
vouloit adopter, même en dépit des équations d'Euler et

\note{le numérateur de la fraction, lu « $\omega c \sqrt{2gb}$ », est d'une
lecture douteuse ; $\zeta$ est un zêta en forme de 3 ; « $LF$ » est écrit
par-dessus un « $EF$ »}

\folio{53r}
de ma note premiere.

Il suffit pour s'en convaincre de remarquer que puisque $\sin u\omega$ est
zéro lorsqu'on fait $u = \frac{1}{2}$, les deux suppositions \struck{\ill{}}
$u > \frac{1}{2}$, $u < \frac{1}{2}$ donneront deux résultats de signes
contraires pour $\sin u\omega$, desorte qu'avant et après le point $L$, les
ordonnées seront de signes différens.

Voilà donc la difficulté entierement résolue pour ce point, Elle venoit de
l'erreur de signe qu'a faite Euler dans l'équation $\gamma = -\gamma'$

Je dois aussi ajouter, contre l'opinion que j'avois avancée dans ma 1ère
note que le facteur $\sin \frac{1}{2}\omega$, donne la solution admissible
$\sin \frac{1}{2}\omega = 0$, ou $\frac{1}{2}\omega = k\pi$. Quant à l'autre
solution, contre laquelle je ne vois pas d'objection, il me semble qu'on ne
peut la rejetter par cette seule raison que les sons rendus par la lame, dans
ses deux portions, ne s'accorderoient pas entr'eux. Les oscillations peuvent
très bien avoir lieu sans être harmoniques.

J'ai cru, Mademoiselle, ne pas devoir vous faire attendre jusqu'à lundi ces
explications que votre discernement appréciera à leur valeur. Je vous les
envoye comme une preuve de mon zèle et de mon dévouement

L. G.

Paris ce 19 Janvier \uncertain{1814}.

\note{le dernier chiffre du millésime se lit 4, mais l'adresse au verso de
ce feuillet (f. 53v) porte un timbre de janvier 1811, la note à laquelle
cette lettre renvoie a été postée en janvier 1811 (f. 43v), et 1811 est
l'année du premier concours ; la lecture « 1811 » est donc possible et la
question reste ouverte. Deux mots biffés avant « $u > \frac{1}{2}$ »}

\page{59}\folio{53v}
A Mademoiselle

Mademoiselle Sophie Germain,

Maison de M. Degligny,

Rue Ste Croix de la bretonnerie

Paris

\note{adresse de la lettre du 19 janvier, au verso de son second feuillet,
écrite en travers ; timbres de la poste « Janv. » avec un quantième dans le
cercle intérieur et « 1811 », marque « G 5 » ; le texte de la lettre
transparoît}

\section{Lettre de Legendre, 31 décembre 1819}

\folio{54r}
Votre équation, Mademoiselle, est très juste; mais la conséquence que vous
en tirez n'est pas admissible; car de ce que $p^5 = A(p^2 - q)$, il ne
s'ensuit pas que $q$ est divisible par $p^2$; on a exactement
$\frac{q}{p^2} = 1 - \frac{p^3}{A}$; mais $\frac{p^3}{A}$ est-il un entier,
parceque $\frac{p^5}{A}$ en est un? Voilà le nœud

Je vous préviens au reste que depuis que je vous ai parlé pour la premiere
fois de ce moyen de recherche, l'opinion que j'avois \struck{qu'il} qu'il
pouvoit réussir est maintenant bien affoiblie, \struck{\ill{}} et qu'en
somme je crois qu'il sera aussi stérile que bien d'autres. C'est pourquoi
vous ferez très bien de ne pas vous en occuper davantage, de peur de perdre
un tems qui peut être employé beaucoup plus utilement à d'autres recherches.

Recevez, mademoiselle, mes hommages respectueux et tous les souhaits qu'une
véritable affection peut inspirer

Le Gendre

31 déc. 1819.

\note{un mot biffé en début de ligne avant « et qu'en somme » ; la page en
regard (f. 53v) est l'adresse transcrite ci-dessus ; f. 54v et f. 55 sont
blancs (vue 60), cette lettre n'ayant pas conservé d'adresse}

\end{document}
